% history-id: equations
\section{Równania ruchu}

Warunki początkowe zapisujemy w postaci
\[
\vect{r}(0)=(0,0),
\qquad
\vect{v}(0)=
\left(v_0\cos\alpha,\;v_0\sin\alpha\right).
\]
Przyspieszenie jest stałe:
\[
\vect{a}(t)=(0,-g).
\]

Ponieważ w kierunku poziomym nie działa przyspieszenie, składowa pozioma
prędkości jest stała:
\[
v_x(t)=v_0\cos\alpha.
\]

W kierunku pionowym prędkość zmienia się liniowo:
\[
v_y(t)=v_0\sin\alpha-gt.
\]

Po scałkowaniu otrzymujemy położenie
\[
x(t)=v_0\cos\alpha\,t,
\]
\[
y(t)=v_0\sin\alpha\,t-\frac{1}{2}gt^2.
\]

Cały ruch można więc zapisać jednym równaniem wektorowym:
\[
\vect{r}(t)=
\left(
v_0\cos\alpha\,t,\;
v_0\sin\alpha\,t-\frac{1}{2}gt^2
\right).
\]
